On veut préparer une solution contenant $\qty{0.10}{\mol\per\L}$ en ions
\ch{Al^{3+}_{(aq)}} à l'aide d'une fiole jaugée de $\qty{100}{\mL}$.

Pour cela on pèse une masse $m$ de sulfate d'aluminium anhydre (\ch{Al2(SO4)3})
que l' on introduit dans la fiole jaugée et que l'on complète avec de l'eau
distillée jusqu'au trait de jauge.
\begin{enumerate}
    \item Exprimer la concentration molaire $c$ du soluté (sulfate d'aluminium)
          en fonction de $m$, $V$ et $M(\ch{Al2(SO4)3})$.
    \item Écrire l'équation de dissolution du soluté dans l'eau.
    \item Pourquoi chaque ion s'entoure t-il de molécule d'eau~? Comment appelle
          t-on ce phénomène ? Prenez l'exemple de l'ion \ch{Al^{3+}_{(aq)}} et
          dessinez les molécules d'eau autour de cet ion.
    \item Calculer la masse $m$ de soluté à peser (Il faut penser à la relation
          entre concentration en soluté et concentration en
          ions).
\end{enumerate}

\begin{donnees}
    en $\unit{\g\per\mol}$~: $M(\ch{Al})=27,0$~; $M(\ch{S})=32,1$~;
    $M(\ch{O})=16,0$.
\end{donnees}
